% Original fixed-tier evaluation fixture for Set P (not a real paper).
% Tests position-sensitive references and an inner equation used in a deferred proof.
\documentclass{article}
\usepackage{amsmath,amssymb,amsthm}

\newtheorem{theorem}{Theorem}[section]
\newtheorem{lemma}[theorem]{Lemma}
\newtheorem{proposition}[theorem]{Proposition}
\newtheorem{corollary}[theorem]{Corollary}
\theoremstyle{definition}
\newtheorem{definition}[theorem]{Definition}
\theoremstyle{remark}
\newtheorem{remark}[theorem]{Remark}

\begin{document}

\section{Introduction}
\label{sec:intro}
We study exponential energy decay for linear dissipative evolutions generated
by coercive symmetric bilinear forms on a real Hilbert space. Our main
estimate is Theorem~\ref{thm:decay}, which upgrades monotonicity of the
squared norm to an explicit exponential decay rate.

\section{Coercive forms and energy identity}
\label{sec:prelim}

\begin{definition}
\label{def:coercive}
Let $H$ be a real Hilbert space with inner product $\langle \cdot, \cdot \rangle$
and norm $\|\cdot\|$. A bounded symmetric bilinear form $a : H \times H \to \mathbb{R}$
is called \emph{$\alpha$-coercive} (for $\alpha > 0$) if
$a(u,u) \ge \alpha \|u\|^2$ for all $u \in H$.
\end{definition}

\begin{lemma}
\label{lem:energy}
Let $a : H \times H \to \mathbb{R}$ be a bounded symmetric bilinear form with
associated bounded self-adjoint operator $A : H \to H$ satisfying
$\langle A u, v \rangle = a(u,v)$ for all $u,v \in H$. If
$u \in C^1([0,\infty); H)$ solves $u'(t) + A u(t) = 0$, then for every $t \ge 0$,
\begin{equation}
\label{eq:energy-id}
\frac{d}{dt}\|u(t)\|^2 + 2 a(u(t),u(t)) = 0.
\end{equation}
\end{lemma}
\begin{proof}
Differentiating $\|u(t)\|^2 = \langle u(t), u(t) \rangle$ and using symmetry
of the inner product gives
\[
\frac{d}{dt}\|u(t)\|^2 = 2 \langle u'(t), u(t) \rangle
= -2 \langle A u(t), u(t) \rangle = -2 a(u(t), u(t)).
\]
Rearranging yields the claimed identity.
\end{proof}

\section{Main decay results}
\label{sec:main}

\begin{proposition}
\label{prop:contraction}
Suppose $a$ is $\alpha$-coercive in the sense of Definition~\ref{def:coercive}.
Then every solution $u \in C^1([0,\infty); H)$ of $u'(t) + A u(t) = 0$ satisfies
$\|u(t)\| \le \|u(0)\|$ for all $t \ge 0$.
\end{proposition}
\begin{proof}
Because $a(u(t), u(t)) \ge \alpha \|u(t)\|^2 \ge 0$, Lemma~\ref{lem:energy}
implies $\frac{d}{dt}\|u(t)\|^2 = -2 a(u(t), u(t)) \le 0$ for all $t \ge 0$.
Integrating from $0$ to $t$ yields $\|u(t)\|^2 \le \|u(0)\|^2$, and taking
square roots concludes the argument.
\end{proof}

\begin{theorem}
\label{thm:decay}
Suppose $a$ is $\alpha$-coercive in the sense of Definition~\ref{def:coercive}
with $\alpha > 0$. Then every solution $u \in C^1([0,\infty); H)$ of
$u'(t) + A u(t) = 0$ satisfies
\[
\|u(t)\| \le e^{-\alpha t} \|u(0)\| \qquad \text{for all } t \ge 0.
\]
\end{theorem}

In contrast to the differential identity of Lemma~\ref{lem:energy}, the global bound requires coercivity.

\begin{remark}
\label{rem:sharpness}
The squared-norm decay rate $2\alpha$ is sharp: if $v \in H \setminus \{0\}$
satisfies $A v = \alpha v$, then $u(t) = e^{-\alpha t} v$ solves
$u'(t) + A u(t) = 0$ and attains $\|u(t)\|^2 = e^{-2\alpha t}\|u(0)\|^2$ for
every $t \ge 0$.
\end{remark}

\begin{corollary}
\label{cor:uniqueness}
Under the hypotheses of Theorem~\ref{thm:decay}, any two solutions
$u_1, u_2 \in C^1([0,\infty); H)$ of $u'(t) + A u(t) = 0$ with
$u_1(0) = u_2(0)$ coincide for all $t \ge 0$.
\end{corollary}
\begin{proof}
Set $w(t) = u_1(t) - u_2(t)$. By linearity, $w$ solves $w'(t) + A w(t) = 0$
with $w(0) = 0$. Applying Proposition~\ref{prop:contraction} gives
$\|w(t)\| \le \|w(0)\| = 0$ for all $t \ge 0$, so $u_1(t) = u_2(t)$.
\end{proof}

\section{Proof of the decay theorem}
\label{sec:deferred-proofs}

\begin{proof}[Proof of Theorem~\ref{thm:decay}]
Set $E(t) = \|u(t)\|^2$. By Proposition~\ref{prop:contraction}, $E$ is
nonincreasing and remains bounded by $E(0)$ on $[0,\infty)$. Combining the
coercivity bound $a(u(t), u(t)) \ge \alpha \|u(t)\|^2$ with
\eqref{eq:energy-id} gives
\[
E'(t) + 2\alpha E(t) \le 0 \qquad \text{for all } t \ge 0.
\]
Multiplying by the integrating factor $e^{2\alpha t}$ shows that
$\frac{d}{dt}\bigl(e^{2\alpha t} E(t)\bigr) \le 0$. Integrating on $[0,t]$
yields $E(t) \le e^{-2\alpha t} E(0)$, and taking square roots completes the
proof.
\end{proof}

\end{document}
